\section{Equivalence Contract}

The contract is serial equivalence, not approximate agreement. For a package in
the current implemented surface, the two verifier entry points must satisfy:

\begin{center}
\small
\begin{tabularx}{\textwidth}{lX}
\toprule
Property & Meaning \\
\midrule
Same pass records & The \texttt{passed} vector contains the same
\texttt{EquationPass} records. \\
Same order & Those records appear in the same order that
\texttt{verify\_package} would have emitted. \\
Same failure behavior & A package rejected by the parallelized check family is
rejected with the same \texttt{RcountCoreError} as the serial verifier. \\
Same remaining surface & Checks not yet parallelized are appended by the
parallel entry point in the same sequence as the serial entry point. \\
\bottomrule
\end{tabularx}
\end{center}

In code, an \texttt{EquationPass} contains an equation id, contest id, and
reporting-unit id. Transcript equality therefore means more than both verifiers
returning success. It means the public report names the same equations for the
same contests and reporting units in the same deterministic sequence.

The selection-sum checks may execute concurrently, but their reports are
collected into a vector before the rest of the verifier surface is appended.
That vector is the boundary where parallel work is converted back into the
serial transcript order. The remaining checks continue to run through the same
serial functions, so the first parallel slice is constrained to be a drop-in
replacement for the opening selection-sum loop.
